\documentclass[11pt,a4paper]{article}
\usepackage[margin=1in]{geometry}
\usepackage[T1]{fontenc}
\usepackage[utf8]{inputenc}
\usepackage{lmodern}
\usepackage{microtype}
\usepackage{enumitem}
\usepackage{xcolor}
\usepackage{hyperref}
\hypersetup{colorlinks=true, linkcolor=black, urlcolor=blue}
\setlength{\parindent}{0pt}
\setlength{\parskip}{0.55em}
\setlist[itemize]{leftmargin=1.5em,itemsep=0.2em,topsep=0.2em}
\setlist[enumerate]{leftmargin=1.7em,itemsep=0.2em,topsep=0.2em}
\pagenumbering{arabic}
\begin{document}
\begin{center}
{\LARGE\bfseries WhatsApp Automated Reply System Proposal}
\end{center}
\vspace{0.5em}

Hello,\par

My name is Fabricio Santana. I am a software engineer and technical lead with 20+ years of experience in backend development, software architecture, database integration, automated testing, cloud computing, and delivery of critical systems. I have also led large engineering teams and worked on projects where data consistency, business rules, validation, security, and maintainability are essential.\par

Your project is a strong match for my background because the main challenge is not only sending automatic replies. The important part is defining a compliant and reliable customer-service workflow where triggers, fallback behavior, routing, Google Sheets integration, and human escalation work together as a maintainable operational system.\par

My relevant approach for this kind of project is to combine backend API integration, automation rules, support workflow design, WhatsApp-compliant provider selection, data integration, logging, deployment support, and automated testing into one coherent delivery plan. This keeps the first version practical while preserving a solid technical base for future expansion.\par

\section*{Working methodology}

My delivery method is designed to reduce risk before development starts. The work begins with a fixed Discovery Sprint and then moves into a delivery package only after scope, priorities, responsibilities, dependencies, and acceptance criteria are clear.\par

The first step is a one-week Discovery Sprint with a fixed price of USD 500. In this phase, I can hold up to five remote meetings to understand the business problem, review the current platform, analyze the current data sources and workflows, map requirements, identify risks, review integrations, and define priorities. The deliverable is a practical work plan for the implementation, including recommended scope, technical approach, milestones, success criteria, and the most appropriate delivery package.\par

After discovery, the project can move into one of three fixed-price, four-week delivery packages. The package is selected according to scope, complexity, delivery speed, integration needs, quality requirements, and the amount of engineering support required. Some roles may participate part-time depending on the agreed scope, especially Scrum Master, technical leadership, test automation, DevOps, observability, and deployment support.\par

\textbf{Delivery package options}\par

\begin{itemize}
\item \textbf{Essential Delivery} --- USD 2,000, 4 weeks. Includes a Scrum Master and one full-stack engineer. This package is designed for a focused module, MVP, backend feature, import flow, or small integration with limited uncertainty. It is suitable when the scope is well defined, the number of integrations is limited, and the main goal is to deliver a reliable working feature without unnecessary overhead.
\item \textbf{Product Acceleration} --- USD 4,000, 4 weeks. Includes a Scrum Master, two full-stack engineers, and one engineer focused on test automation, DevOps, observability, and deployment. This package is recommended when faster delivery is needed or when the scope includes a product, API, integration, internal workflow, multiple features, or more demanding quality needs.
\item \textbf{Scale \& Intelligence} --- USD 6,000, 4 weeks. Includes a Technical lead, Scrum Master, two full-stack engineers, and one engineer focused on test automation, DevOps, observability, and deployment. This package is designed for complex systems, advanced data workflows, automation, auditability, security, observability, broader platform evolution, or stronger technical scale planning.
\end{itemize}

Role responsibilities are clear: the Scrum Master organizes scope, communication, ceremonies, progress tracking, and delivery alignment; full-stack engineers implement backend, frontend, API, and database changes as needed; the test automation, DevOps, observability, and deployment engineer supports automated tests, CI/CD, monitoring, release quality, and deployment; and the Technical lead, included in the Scale \& Intelligence package, owns architecture, technical decisions, code quality, and risk management for more complex projects.\par

For an automation project like this, the implementation setup can include provider integration, routing logic, Sheets synchronization, logging, operational monitoring, deployment, and handoff work within the agreed package scope.\par

Each delivery cycle is executed with Scrum-inspired practices, frequent alignment, technical review, automated testing when applicable, and a demonstration of what was delivered. The client keeps ownership of the produced code, and the scope is agreed before execution to avoid open-ended work and unclear expectations.\par

\textbf{Construction defect warranty.} After delivery and acceptance, I include a 30-day warranty for construction defects related to the approved scope. This covers fixes for implementation errors, broken agreed behavior, or defects introduced by the delivered code. It does not cover new features, changes in business rules, new message templates or routing rules not included in the approved scope, third-party service changes, infrastructure changes outside the project, or requests that expand the original acceptance criteria.\par

\section*{Opportunity analysis}

This project should be treated as a customer-service workflow automation problem, not as a simple chatbot setup. A reliable WhatsApp auto-responder needs more than keyword replies. It needs a compliant provider choice, clear session and template behavior, reliable rule execution, Sheets integration, conversation logging, handoff to human agents, and operational visibility after go-live.\par

The business value is also important: the system should reduce response time, standardize answers to frequent questions, reduce repetitive manual work, improve handoff to the right agent, and create better visibility into customer interactions.\par

The main technical risks are selecting the right WhatsApp-compliant integration model, defining maintainable routing and escalation rules, choosing the right role for Google Sheets, keeping the first version bounded, and avoiding expansion into CRM, AI assistants, dashboards, or broader multichannel support before the core workflow is stable. I recommend avoiding unofficial WhatsApp Web automation because it can create account-ban, reliability, and maintenance risks. For this reason, I recommend starting with a controlled and explicit rule-based workflow rather than an open-ended chatbot design.\par

A practical first implementation cycle can start with a compliant WhatsApp integration, FAQ auto-replies, explicit escalation rules, Sheets-based lookup or logging according to the agreed model, operational logs, and the technical foundation required for future expansion. The exact delivery boundary should be confirmed during discovery.\par

\section*{Proposed work plan}

I would approach the project in structured phases. The Discovery Sprint is the immediate recommended commitment, and implementation can move directly into a fixed delivery cycle once the provider, workflow, and acceptance criteria are confirmed.\par

\textbf{1. Discovery, workflow definition, and technical design}\par

\begin{itemize}
\item Review the current support process, message types, FAQ patterns, escalation rules, support-team workflow, and expected message volume.
\item Define the preferred WhatsApp-compliant provider and the correct integration approach for the business case.
\item Identify how Google Sheets should be used in the first version: lookup source, logging layer, routing table, or operational support table.
\item Produce a concrete implementation plan with workflow map, recommended architecture, acceptance criteria, responsibilities, dependencies, risks, and package recommendation.
\end{itemize}

\textbf{2. Core integration architecture and backend foundation}\par

\begin{itemize}
\item Design and implement the integration foundation for incoming messages, reply rules, Sheets access, routing logic, and logging.
\item Define data structures and APIs that preserve traceability and allow future extension.
\item Add validation, authorization, auditability, and structured error handling.
\item Establish the technical base for deployment, monitoring, and secure environment configuration.
\end{itemize}

\textbf{3. Automated reply and FAQ workflow}\par

\begin{itemize}
\item Implement the first set of predefined triggers, keywords, FAQ responses, and fallback behavior.
\item Keep the first implementation cycle focused on a bounded support workflow rather than feature overload.
\item Structure the response logic so it can be maintained and expanded safely after go-live.
\item Support dynamic content usage when Google Sheets is part of the agreed reply model.
\end{itemize}

\textbf{4. Human escalation and operational routing}\par

\begin{itemize}
\item Implement explicit escalation rules for cases the automation cannot resolve.
\item Route conversations to the correct human agent or support flow according to the agreed business rules.
\item Preserve logs and decision traceability for routed interactions.
\item Define fallback behavior for integration failures, missing data, or unsupported message paths.
\end{itemize}

\textbf{5. Testing, monitoring, deployment support, and handoff}\par

\begin{itemize}
\item Test message handling, reply rules, escalation flows, Sheets integration, and edge cases.
\item Add operational visibility for failed replies, failed Sheets calls, routing issues, and unusual usage patterns.
\item Document the architecture, provider setup, workflow limitations, and maintenance steps.
\item Provide deployment support and a handoff session or written notes so your team can operate and extend the system safely.
\end{itemize}

\section*{Estimated timeline}

For the immediate next step, my recommendation is a 1-week Discovery Sprint. After discovery, the first implementation cycle can move into a fixed 4-week delivery package. For a bounded WhatsApp automation workflow with FAQ replies, routing, and Google Sheets integration, the likely recommendation is either one Essential Delivery cycle or one Product Acceleration cycle, depending on the operational complexity confirmed during discovery.\par

A practical schedule would be:\par

\begin{itemize}
\item Week 1: Discovery Sprint, workflow mapping, provider recommendation, Sheets usage model, launch-scope backlog, acceptance criteria, and implementation plan.
\item Weeks 2 to 5: first fixed implementation cycle, covering the agreed automation workflow and operational handoff.
\end{itemize}

If discovery shows that the first version is narrower than currently implied, the implementation may fit the smaller package cleanly. If it shows higher routing complexity, stronger observability needs, broader support flows, or additional integrations, I will recommend the larger package or additional cycles accordingly.\par

\section*{Availability and communication}

I can work remotely and keep communication in English or Portuguese through Workana messages, scheduled calls, and written progress updates. I usually prefer short alignment checkpoints during the week, especially after reviewing the workflow and after each major implementation milestone.\par

This helps reduce misunderstandings and keeps the project aligned with your actual support process, business rules, and operational workflow.\par

\section*{Budget}

Based on the current understanding, my recommended immediate starting budget is USD 500 for the one-week Discovery Sprint. This gives you a concrete plan before committing to the full implementation budget, while keeping a clear path to delivery in the following implementation cycle.\par

After discovery, the first implementation cycle will likely fit one of these packages:\par

\begin{itemize}
\item \textbf{Essential Delivery} at USD 2,000 for 4 weeks if the scope is a bounded FAQ and routing workflow with straightforward Sheets integration.
\item \textbf{Product Acceleration} at USD 4,000 for 4 weeks if the scope includes richer routing rules, stronger monitoring, broader operational handoff, or more advanced business logic.
\end{itemize}

The Discovery Sprint will refine the scope and confirm the correct implementation package. If the review shows a smaller and tightly bounded workflow, I can recommend the lower package. If it shows broader support operations, higher message volume, stronger reliability needs, or additional integrations, I will recommend the larger package or additional cycles.\par

\section*{Questions before final confirmation}

Before confirming the final scope, I would like to clarify a few points:\par

\begin{enumerate}
\item Which WhatsApp provider or business setup is currently available: Meta official API, Twilio, Zenvia, 360dialog, or another provider?
\item Should the first version be fully rule-based, or do you want AI-assisted intent handling later?
\item How should Google Sheets be used in practice: FAQ source, interaction log, routing table, dynamic reply source, or more than one of these?
\item What should trigger escalation to a human agent?
\item How many agents will receive routed conversations, and how are they assigned today?
\item What message volume do you expect per day or per month?
\item Are approved message templates already configured, or should this be included in the setup plan?
\item Should the system only reply and route, or also maintain status, audit logs, and operational metrics?
\item Is there already an approved budget range and a target date for the first operational version?
\item Who will approve access to the WhatsApp provider, Google account, and deployment environment?
\end{enumerate}

Once these points are clear, I can confirm the final delivery plan and milestones.\par

\end{document}
